\documentclass[12pt]{article}
\usepackage{graphicx} % Required for inserting images
\usepackage[margin=0.5in]{geometry}
\usepackage{amsmath, amssymb}
\usepackage{mathrsfs}


\usepackage{soul}


% Dr. David Brown Commands
\newcommand{\ds}{\displaystyle}
\newcommand{\R}{\mathbb{R}}
\newcommand{\N}{\mathbb{N}}
\newcommand{\Q}{\mathbb{Q}}
\newcommand{\C}{\mathbb{C}}


% Commands to get script r
\usepackage{calligra}
\DeclareMathAlphabet{\mathcalligra}{T1}{calligra}{m}{n}
\DeclareFontShape{T1}{calligra}{m}{n}{<->s*[2.2]callig15}{}
\newcommand{\scripty}[1]{\ensuremath{\mathcalligra{#1}}}

% Unit commands
\newcommand{\degree}{\text{°}}
\newcommand{\unitSpeed}{\frac{\text{m}}{\text{s}}}
\newcommand{\unitAcceleration}{\frac{\text{m}}{\text{s}^2}}

% Constant commands
\newcommand{\avogadro}{6.022\times10^{23}}

\newcommand{\boltzmannConstK}{1.381\times10^{-23}\frac{\text{J}}{\text{K}}}
\newcommand{\planckConst}{6.626\times10^{-34}\text{J}\cdot\text{s}}

\newcommand{\atmP}{1.01325\times10^5\,\text{Pa}}

% Commands to easily write partial derivatives
\newcommand{\partDeriv}[2]{\frac{\partial #1}{\partial #2}}
\newcommand{\doublePartDeriv}[2]{\frac{\partial^2 #1}{\partial^2 #2}}


\title{Problem 1.16}
\author{Gabriel Decker}


\begin{document}



\maketitle
\begin{flushleft}

First, let's consider the forces on the horizontal slab of air.
There are the forces from the weight above the slab, the force from the slab's weight, and the force pushing up from the air below the slab.
The force from the slab's weight are pretty well understood.
The others are not.
Label the force from above $F_a$ and the force from below $F_b$.
These are in the z-direction and point downward and upward, respectively.\\~\\

Since the forces must cancel out for the system to be in equilibrium, this implies that
$$-F_a-mg+F_b=0$$
where here $m$ is the total mass of the slab.
If our slab of air has a surface area $A$ and thickness $dz$, then the mass of the slab can be replaced with the following, given some density $\rho=m/Adz$.
$$\implies -F_a-\rho Agdz+F_b=0$$
Dividing the equation by $A$, we can define the equation in terms of the pressure from the top and bottom instead.
$$\implies -P_a-\rho gdz+P_b=0$$\\~\\

Since our slice has infinitesimal thickness, the difference between the pressure from the top and bottom should be infinitesimal.
Therefore we can define $dP=P_a-P_b$ since increasing altitude is in the positive z direction.
$$\implies dP=-\rho gdz$$
Therefore, 
$$\implies \frac{dP}{dz}=-\rho g$$\\~\\

We must now get an expression of $\rho$ in terms of $T$, $P$ and the average mass of air $m$.
We're going to now say that $m_T$ is the total mass of the slab.
Recall that in general, $\rho=m/V$.
By the ideal gas law, $V=\frac{NkT}{P}$, so this implies that in our case, $\rho=\frac{m_T}{NkT}P$.
We have a problem with the $N$ being in there.
The problem states that we want an average mass $m$.
Say we have $N$ air molecules and a total mass of $m_T$.
The average mass $m$ would then be $m=m_T/N$.
Substituting this into our equation for $\rho$, we get
$$\rho=\frac{m}{kT}P$$
Plugging in this equation for $\rho$, we get the following differential equation.
$$\implies \frac{dP}{dz}=-\frac{mg}{kT}P$$\\~\\

Finally, I'll solve the differential equation.
This one is pretty simple.
It's clearly an exponential with $z$ as the variable because differentiating $P$ with respect to $z$ yields $P$ again with some constants in front. So,
$$\implies P=P_0e^{-mgz/kt}$$\\~\\

But, what is $P_0$?
That is an arbitrary constant, but we can define it, based on some things we know.
We know that at sea level ($z=0$), there is 1 atm of pressure, so define $P_0=P(0)$ as 1 atm (in whatever unit you like).
$$\implies P=P(0)e^{-mgz/kt}$$\\~\\




\end{flushleft}

\end{document}
